\chapter{L'apparato urinario}

% ---------------------------------------------------------------------------
\section{Generalità}

\textbf{Sistema urinario}: insieme di organi con funzioni di \textbf{filtrazione del sangue} ed
eliminazione all'esterno dei cataboliti. È formato dai \textbf{reni} e dalle \textbf{vie urinarie}
(ureteri, vescica, uretra).
\begin{itemize}
  \item i \textbf{reni} realizzano, tramite i \textbf{corpuscoli renali} (di Malpighi), la
    filtrazione del sangue $\rightarrow$ \textbf{pre-urina};
  \item la pre-urina scorre nel sistema tubulare, che riassorbe metaboliti, $H_2O$ e sali secondo le
    necessità $\rightarrow$ \textbf{urina definitiva};
  \item l'urina è trasportata all'esterno dalle vie urinarie, dai reni (papille renali) fino al meato
    uretrale esterno.
\end{itemize}
Il sistema urinario è la ``bilancia del plasma'': ogni rene depura il sangue che vi passa.

\rubrica{Funzioni del rene}
\begin{enumerate}
  \item escrezione dei prodotti terminali del metabolismo (scorie azotate di origine proteica);
  \item controllo dell'equilibrio idrico;
  \item controllo dell'equilibrio degli elettroliti;
  \item controllo dell'equilibrio acido-base;
  \item regolazione della pressione arteriosa (sistema \textbf{renina-angiotensina});
  \item regolazione dell'emopoiesi (secrezione di \textbf{eritropoietina});
  \item regolazione dell'assorbimento intestinale del calcio.
\end{enumerate}

% ---------------------------------------------------------------------------
\section{Il rene}

\textbf{Rene}: organo pari, situato nella porzione posteriore dell'addome (regione lombare), nella
parte superiore dello spazio \textbf{retroperitoneale} (logge renali).
\begin{itemize}
  \item dimensioni: lunghezza 12 (10--14) cm, larghezza 6 (5--7) cm, spessore $\sim$3 cm; peso
    150--160 g; colorazione rosso-bruna;
  \item posizione: da T12 a L3, a lato della colonna; il \textbf{rene dx} è più basso del sx di
    $\sim$2 cm (rapporto con il fegato).
\end{itemize}

\rubrica{Morfologia}
\begin{itemize}
  \item 2 \textbf{poli} (superiore e inferiore, in rapporto con le surrenali);
  \item 2 \textbf{margini}: mediale (concavo) e laterale (convesso);
  \item 2 \textbf{facce}: antero-laterale e postero-mediale;
  \item la superficie concava (verso la colonna) accoglie l'\textbf{ilo renale}: ingresso
    dell'arteria renale e dei nervi, uscita di vena renale, uretere, vasi linfatici.
\end{itemize}

\rubrica{Vascolarizzazione}
\begin{itemize}
  \item i reni ricevono $\sim$22\% (range 20--25\%) della gittata cardiaca ($\sim$600
    cm\textsuperscript{3}/minuto per rene) $\rightarrow$ arteria e vena renale di ampio calibro;
  \item sistema vascolare di tipo \textbf{terminale}; deflusso nelle vene renali $\rightarrow$ vena
    cava inferiore;
  \item sequenza: arteria renale $\rightarrow$ arterie segmentali $\rightarrow$ interlobari
    $\rightarrow$ arcuate $\rightarrow$ corticali radiali (interlobulari) $\rightarrow$ arteriole
    afferenti $\rightarrow$ glomeruli $\rightarrow$ arteriole efferenti $\rightarrow$ arteriole rette
    $\rightarrow$ reti capillari peritubulari $\rightarrow$ venule rette $\rightarrow$ vene corticali
    radiali $\rightarrow$ arcuate $\rightarrow$ interlobari $\rightarrow$ vena renale.
\end{itemize}

% ---------------------------------------------------------------------------
\section{Struttura interna del rene}

Il parenchima, rivestito da una capsula connettivale, si divide in \textbf{corticale} e
\textbf{midollare}.
\begin{itemize}
  \item \textbf{corticale renale}: più chiara e periferica; contiene i \textbf{corpuscoli renali};
    invia le \textbf{colonne renali} verso il seno renale (labirinto corticale + raggi midollari);
  \item \textbf{midollare renale}: più scura e centrale; forma le \textbf{piramidi renali} (di
    Malpighi), separate dalle colonne renali.
\end{itemize}

\rubrica{Piramidi, colonne e lobo}
\begin{itemize}
  \item \textbf{piramidi del Malpighi}: base convessa (appoggiata alla corticale) e apice (papilla)
    che sporge nel seno renale $\rightarrow$ l'apice forma le papille renali che confluiscono nei
    calici minori;
  \item \textbf{colonne renali}: tra le piramidi, formate da corticale;
  \item \textbf{lobo renale}: una piramide + la corticale che la circonda.
\end{itemize}

% ---------------------------------------------------------------------------
\section{Il nefrone}

\textbf{Nefrone}: unità morfofunzionale del rene ($\sim$1 milione per rene; lunghezza totale del
sistema dei nefroni nei 2 reni $>$ 50 km). Comprende una parte filtrante (\textbf{corpuscolo renale})
e una riassorbente-secernente (\textbf{tubulo renale}).
\begin{itemize}
  \item \textbf{glomerulo}: cellule capaci di filtrare il plasma;
  \item \textbf{tubulo renale}: prossimale $\rightarrow$ ansa di Henle $\rightarrow$ distale; i
    tubuli distali confluiscono nel \textbf{tubulo collettore}, più tubuli nel \textbf{dotto
    collettore} $\rightarrow$ immette l'urina nell'uretere.
\end{itemize}
Epitelio: pavimentoso semplice nella capsula glomerulare e nella parte sottile dell'ansa; cubico
semplice nei tubuli convoluti e nel dotto collettore. Nel tubulo prossimale le cellule hanno fitti
microvilli (\textbf{orletto a spazzola}).

\subsection{Corpuscolo renale}
Formazioni rotondeggianti nella corticale (sottocapsulari o nelle colonne); responsabili
dell'\textbf{ultrafiltrazione} del plasma; $\sim$1 milione/rene, superficie filtrante totale 1,5--2
m\textsuperscript{2}.
\begin{itemize}
  \item \textbf{capsula glomerulare} (di Bowman): foglietto parietale e viscerale (podociti),
    separati dallo spazio capsulare in cui si raccoglie la pre-urina;
  \item \textbf{glomerulo renale}: gomitolo di capillari arteriosi (rete mirabile) tra arteriola
    afferente ed efferente;
  \item \textbf{mesangio} intraglomerulare: insieme di \textbf{cellule mesangiali} $\rightarrow$
    sostegno ai capillari, nucleo denso e forma irregolare, inglobano sostanze che potrebbero
    ostruire la lamina densa, regolano il diametro dei capillari (flusso e filtrazione), ruolo nei
    meccanismi di immunità e infiammazione.
\end{itemize}

% ---------------------------------------------------------------------------
\section{La formazione dell'urina}

\subsection{Filtrazione}
Il sangue è filtrato a livello del glomerulo: acqua, soluti a basso peso molecolare (ioni
inorganici), urea, glucosio e amminoacidi fuoriescono dai capillari $\rightarrow$ spazio capsulare
$\rightarrow$ tubulo prossimale. \textbf{Non} passano le sostanze con peso molecolare $>$ 70.000
(proteine plasmatiche).
\begin{itemize}
  \item riassorbite (sostanze utili): acqua, sodio, potassio, cloro, glucosio, amminoacidi;
  \item escrete con l'urina: urea, acido urico, creatinina, ione ammonio, eventuali sostanze
    tossiche.
\end{itemize}

\subsection{Riassorbimento}
Funzione principale del \textbf{tubulo contorto prossimale}: le sue cellule riassorbono con
meccanismo attivo ioni, sostanze organiche e proteine, cedute allo spazio interstiziale dalle reti
capillari peritubulari. Qui si riassorbe oltre il 60\% dell'ultrafiltrato (amminoacidi, glucosio,
cationi $Na^+$/$K^+$/$Ca^{2+}$, anioni come solfati e fosfati).

\begin{table}[htbp]
  \centering \small \renewcommand{\arraystretch}{1.3}
  \begin{tabularx}{\textwidth}{|l|X|}
    \hline
    \textbf{Segmento} & \textbf{Principali riassorbimenti} \\
    \hline
    Tubulo contorto prossimale
      & $Na^+$ 65\%, $Cl^-$ 65\%, acqua 65\%, glucosio 100\%, amminoacidi 100\%, bicarbonato 80\%,
        fosfato 80\%, $Ca$ 70\%, urea 80\%, $K^+$ 70\% \\
    \hline
    Ansa di Henle (tratto ascendente spesso)
      & $Na^+$ 25\%, $Cl^-$ 25\%, bicarbonato 15\%, $Ca$ 10\%, $K^+$ 25\% \\
    \hline
    Tubulo contorto distale
      & $Na^+$/$Cl^-$/$K^+$ 5\% (aldosterone-dip.), $Ca$ (PTH-dip.), protoni e bicarbonato (pH-dip.) \\
    \hline
    Dotto collettore
      & $Na^+$/$Cl^-$ 5\%, acqua 20\% (ADH-dip.), urea 80\% (ADH-dip.), $K^+$ 10\% \\
    \hline
  \end{tabularx}
\end{table}
Legenda: \textit{ADH-dip.} = dipendente dall'ormone antidiuretico; \textit{Ald-dip.} = da
aldosterone; \textit{pH-dip.} = dal pH; \textit{PTH} = paratormone.

\subsection{Sistema RAAS (renina-angiotensina-aldosterone)}
Le cellule della \textbf{macula densa} sono sensibili alla concentrazione di $Na^+$ nel tubulo.
\begin{itemize}
  \item una riduzione della pressione arteriosa $\rightarrow$ minore ultrafiltrato $\rightarrow$
    $\downarrow$ $Na^+$ nei tubuli $\rightarrow$ la macula densa stimola la produzione di
    \textbf{renina} dalle cellule iuxtaglomerulari;
  \item sequenza: renina + angiotensinogeno $\rightarrow$ \textbf{angiotensina I} $\rightarrow$
    (ACE) \textbf{angiotensina II} $\rightarrow$ \textbf{aldosterone} $\rightarrow$ $\uparrow$
    riassorbimento di $Na^+$ e acqua, $\uparrow$ pressione.
\end{itemize}

% ---------------------------------------------------------------------------
\section{Le vie urinarie intrarenali}

\begin{itemize}
  \item \textbf{calici minori}: origine delle vie urinarie; forma conica, 1 per papilla renale;
    raccolgono l'urina sgocciolata dalle papille e si riuniscono in gruppi di 3;
  \item \textbf{calici maggiori}: confluenza dei calici minori (numero e dimensioni variabili);
    ramificazioni della pelvi;
  \item \textbf{pelvi renale} (bacinetto): dalla confluenza dei calici maggiori; forma conica, una
    porzione nel seno renale e una che continua nell'\textbf{uretere}.
\end{itemize}
Morfologia microscopica: tonaca mucosa = epitelio di transizione (\textbf{urotelio}); tonaca
muscolare = fasci di cellule muscolari lisce; tonaca avventizia = connettivo lasso.

% ---------------------------------------------------------------------------
\section{Gli ureteri}

\textbf{Ureteri}: organi cavi che collegano la pelvi renale (da cui originano) alla vescica (dove
terminano con l'orifizio ureterale). Ruolo attivo nel trasporto: la distensione da urina stimola la
muscolatura liscia a contrarsi (peristalsi); forza e frequenza sono correlate all'urina in ingresso.
\begin{itemize}
  \item \textbf{parti}: addominale (retroperitoneale, concavità anteriore), pelvica
    (sottoperitoneale, concavità posteriore), intramurale (nello spessore della parete vescicale);
  \item \textbf{restringimento} all'origine (1\textsuperscript{o} restringimento);
  \item microscopica: mucosa = urotelio, muscolare = più strati, avventizia = connettivo.
\end{itemize}

% ---------------------------------------------------------------------------
\section{La vescica urinaria}

\textbf{Vescica}: organo cavo delle vie urinarie inferiori, deputato alla raccolta dell'urina; nella
piccola pelvi, posteriormente alla sinfisi pubica.
\begin{itemize}
  \item forma e rapporti variano col riempimento: vuota è appiattita (piccola pelvi), piena si spinge
    in alto (grande pelvi);
  \item base (o fondo), corpo, apice; una faccia anteriore, una posteriore, due laterali; le due
    laterali e la posteriore convergono nel \textbf{collo vescicale}, da cui origina l'uretra;
  \item \textbf{trigono}: triangolo con gli angoli superiori agli sbocchi degli ureteri e l'inferiore
    all'orifizio uretrale interno;
  \item legamenti la connettono agli organi circostanti.
\end{itemize}

\rubrica{Struttura}
\begin{itemize}
  \item mucosa: \textbf{urotelio} impermeabile $\rightarrow$ impedisce all'urina ipertonica di
    richiamare acqua dagli interstizi; distensibile senza cedimenti; assenza di ghiandole;
  \item muscolare: 3 fasci $\rightarrow$ muscolo \textbf{detrusore} e \textbf{sfintere interno}
    dell'uretra;
  \item avventizia (connettivo denso), sostituita dalla sierosa (peritoneo) nelle porzioni superiori.
\end{itemize}

\rubrica{Innervazione}
\begin{itemize}
  \item \textbf{vescica piena}: i recettori di distensione, via parasimpatica (midollo sacrale
    S2--S4), determinano il rilascio dello sfintere interno (liscio) $\rightarrow$ svuotamento;
  \item \textbf{vescica vuota}: prevale il tono simpatico (L1--L3) $\rightarrow$ vescica rilasciata,
    sfintere interno contratto.
\end{itemize}

% ---------------------------------------------------------------------------
\section{L'uretra}

\textbf{Uretra}: organo cavo, impari e mediano, deputato al trasporto dell'urina dalla vescica
all'esterno; decorso e rapporti diversi tra maschio e femmina. Compresa tra due sfinteri: interno
(liscio) ed esterno (striato $\rightarrow$ controllo volontario). Epitelio di transizione, squamoso
cheratinizzato presso l'orifizio esterno.
\begin{itemize}
  \item \textbf{maschile}: $\sim$15--20 cm; doppia funzione (urina + passaggio dello sperma
    nell'eiaculazione); segmenti: \textbf{preprostatica}, \textbf{prostatica}, \textbf{membranosa}
    (attraversa il diaframma urogenitale), \textbf{spongiosa} (nel corpo spongioso del pene, fino
    all'orifizio uretrale esterno);
  \item \textbf{femminile}: più corta ($\sim$4--5 cm), più semplice; attraversa il diaframma
    urogenitale $\rightarrow$ parte pelvica e parte perineale.
\end{itemize}

% ---------------------------------------------------------------------------
\section{Le urine}

In condizioni normali il volume nelle 24 ore è $\sim$1500--2000 mL.
\begin{itemize}
  \item \textbf{poliuria} (aumento): diabete insipido (carenza di ADH); diabete mellito (il glucosio
    nelle urine richiama acqua per osmosi);
  \item \textbf{oliguria} (riduzione): es.\ glomerulonefrite;
  \item \textbf{esame} fisico/chimico/biochimico: colore, aspetto, pH, densità; ricerca di linfociti,
    proteinuria, nitrati, glucosio.
\end{itemize}

% ---------------------------------------------------------------------------
\section{Note cliniche}

\rubrica{Calcolosi urinaria}
Tra i disturbi più dolorosi e comuni delle vie urinarie. I calcoli sono aggregati di sali minerali
dell'urina (ossalato di calcio, fosfato di calcio, acido urico); per lo più espulsi normalmente col
flusso urinario se di dimensioni $<$ 5 mm.

\rubrica{Insufficienza renale cronica}
\begin{itemize}
  \item cause: diabete mellito, ipertensione arteriosa, ostruzione delle vie urinarie, nefriti
    interstiziali;
  \item diagnosi: esami delle urine ed ematochimici; la malattia renale progredisce per stadi fino
    all'insufficienza terminale (dialisi).
\end{itemize}
